\documentclass[11pt]{article}

\usepackage[utf8]{inputenc}
\usepackage[T1]{fontenc}
\usepackage{lmodern}
\usepackage{geometry}
\usepackage{amsmath,amssymb,amsfonts,amsthm,mathtools}
\usepackage{bm}
\usepackage{enumitem}
\usepackage{microtype}
\usepackage{hyperref}
\usepackage{titlesec}
\usepackage{booktabs}
\usepackage{array}
\usepackage{setspace}
\usepackage{mathrsfs}
\usepackage{physics}

\geometry{margin=1in}
\hypersetup{colorlinks=true, linkcolor=black, urlcolor=blue, citecolor=black}
\setlist[enumerate]{topsep=0.4em,itemsep=0.35em}
\setlist[itemize]{topsep=0.3em,itemsep=0.25em}
\titleformat{\section}{\large\bfseries}{\thesection.}{0.5em}{}
\titleformat{\subsection}{\normalsize\bfseries}{\thesubsection.}{0.5em}{}
\setstretch{1.06}

\newcommand{\Aether}{\AE{}ther}
\newcommand{\Aflow}{\AE{}ther-flow}
\newcommand{\TheoryName}{\Aflow{} Interpretation of Relativity}
\newcommand{\TheoryProgram}{\Aether{} / \Aflow{} framework}

\newcommand{\CompletedMetric}{g^{\sharp}}
\newcommand{\MatterSpace}{\Psi}

\newcommand{\Car}{\mathrm{car}}
\newcommand{\Cur}{\mathrm{cur}}
\newcommand{\Hom}{\mathrm{hom}}

\newcommand{\ForSpace}[1]{\mathcal I_{#1}^{\mathrm{for}}}
\newcommand{\PrimShell}{\mathcal S_{\mathrm{prim}}}
\newcommand{\Reduction}[1]{\mathcal R_{#1}}
\newcommand{\CurrentCarrier}[1]{\mathfrak C_{#1}^{\mathrm{cur}}}
\newcommand{\EqCarrier}[1]{\mathfrak C_{#1}^{\mathrm{eq}}}
\newcommand{\MetricOut}{\mathscr G_{\mathrm{out}}}
\newcommand{\MatterOut}{\mathscr T_{\mathrm{out}}}

\newcommand{\SubTarget}[1]{E_{#1}^{\mathrm{sub}}}
\newcommand{\SubPair}[1]{\mathfrak s_{#1}^{\mathrm{sub}}}
\newcommand{\EqnTarget}[1]{Q_{#1}^{\mathrm{eqn}}}
\newcommand{\HiddenSpace}[1]{\mathcal H_{#1}^{\mathrm{hid}}}
\newcommand{\HiddenBody}[1]{D_{#1}^{\mathrm{hid}}}

\newcommand{\ProtoSpace}[1]{\mathcal U_{#1}^{\mathrm{proto}}}
\newcommand{\ProtoBody}[1]{\mathcal B_{#1}^{\mathrm{proto}}}
\newcommand{\ProtoSection}[1]{\Gamma_{#1}^{\mathrm{proto}}}
\newcommand{\ProtoShellSet}[1]{\mathcal Z_{#1}^{\mathrm{proto}}}
\newcommand{\ProtoZeroCore}[1]{\mathcal Z_{#1}^{0,\mathrm{proto}}}
\newcommand{\ProtoSplit}[1]{\Phi_{#1}^{\mathrm{proto}}}
\newcommand{\ProtoCharge}[1]{\mathcal Q_{#1}^{\mathrm{proto}}}
\newcommand{\ProtoEmbed}[1]{\zeta_{#1}^{\mathrm{proto}}}
\newcommand{\ProtoKernelCh}[1]{\widetilde K_{#1}^{0,\mathrm{proto}}}
\newcommand{\ProtoStateComp}[1]{P_{#1}^{\mathrm{proto,co}}}

\newcommand{\PreProtoNucleus}{\mathfrak Q^{\mathrm{nuc}}}
\newcommand{\PreProtoSpace}[1]{\mathcal V_{#1}^{\mathrm{preproto}}}
\newcommand{\PreProtoBody}[1]{\mathcal B_{#1}^{\mathrm{preproto}}}
\newcommand{\PreProtoProj}[1]{\Pi_{#1}^{\mathrm{preproto}}}
\newcommand{\PreProtoProtoMin}[1]{\mathfrak f_{#1}^{\mathrm{preproto,proto}}}
\newcommand{\PreProtoSection}[1]{\Gamma_{#1}^{\mathrm{preproto}}}
\newcommand{\PreProtoFunctional}[1]{\mathscr W_{#1}^{\mathrm{preproto}}}
\newcommand{\PreProtoShellSet}[1]{\mathcal Z_{#1}^{\mathrm{preproto}}}

\newcommand{\PreNucCore}{\mathfrak R^{\mathrm{core}}}
\newcommand{\PreNucSpace}[1]{\mathcal W_{#1}^{\mathrm{prenuc}}}
\newcommand{\PreNucBody}[1]{\mathcal B_{#1}^{\mathrm{prenuc}}}
\newcommand{\PreNucProj}[1]{\Pi_{#1}^{\mathrm{prenuc}}}
\newcommand{\PreNucNucleusMin}[1]{\mathfrak f_{#1}^{\mathrm{prenuc,nuc}}}
\newcommand{\PreNucShellSet}[1]{\mathcal Z_{#1}^{\mathrm{prenuc}}}

\newcommand{\PreCoreTransportCore}{\mathfrak S^{\mathrm{tr}}}
\newcommand{\PreCoreTrBody}[1]{\mathcal B_{#1}^{\mathrm{precore,tr}}}
\newcommand{\PreCoreTrProj}[1]{\widehat{\Pi}_{#1}^{\mathrm{precore}}}
\newcommand{\PreCoreTrLift}[1]{\mathfrak f_{#1}^{\mathrm{precore,tr}}}
\newcommand{\PreCoreTrShell}[1]{\mathcal Z_{#1}^{\mathrm{precore,tr}}}

\newcommand{\SubPreCorePackage}{\mathfrak U^{\mathrm{subpre}}}
\newcommand{\SubPreCoreSpace}[1]{\mathcal L_{#1}^{\mathrm{subpre}}}
\newcommand{\SubPreCoreBody}[1]{\mathcal B_{#1}^{\mathrm{subpre}}}
\newcommand{\SubPreCoreProj}[1]{\Pi_{#1}^{\mathrm{subpre}}}
\newcommand{\SubPreCoreTransportMin}[1]{\mathfrak f_{#1}^{\mathrm{subpre,tr}}}
\newcommand{\SubPreCoreShellSet}[1]{\mathcal Z_{#1}^{\mathrm{subpre}}}

\newcommand{\SubPreImageCore}{\mathfrak U^{\mathrm{img}}}
\newcommand{\SubPreImgBody}[1]{\mathcal B_{#1}^{\mathrm{subpre,img}}}
\newcommand{\SubPreImgProj}[1]{\widehat{\Pi}_{#1}^{\mathrm{subpre}}}
\newcommand{\SubPreImgShell}[1]{\mathcal Z_{#1}^{\mathrm{subpre,img}}}

\newcommand{\SubPreSectionCore}{\mathfrak U^{\mathrm{sec}}}
\newcommand{\SubPreSection}[1]{\mathfrak s_{#1}^{\mathrm{subpre,sec}}}

\newtheorem{definition}{Definition}
\newtheorem{proposition}{Proposition}
\newtheorem{remark}{Remark}
\newtheorem{corollary}{Corollary}
\newtheorem{theorem}{Theorem}

\title{\textbf{The \TheoryName{}}\\[0.4em]
\large Primitive-Reservoir Observer-Reduction\\
\large Observer-Relevant Sub-Pre-Core Exact-Shell Transport-Image Core\\
\large Collapse to an Observer-Relevant Sub-Pre-Core Exact-Shell Transport-Section Core}
\author{}
\date{}

\begin{document}

\maketitle

\begin{abstract}
The current primitive-reservoir observer-reduction frontier already sharpened the live burden inside the exact-shell-transport-core-generating substrate sub-pre-core package: the full sub-pre-core package was compressed to the observer-relevant sub-pre-core exact-shell transport-image core. The next honest move below that theorem is therefore not another same-output deeper-origin relay that leaves the image core unchanged. The sharper benchmark-facing move is to ask whether the image core itself still carries presentational redundancy that can be collapsed away without changing the benchmark one-metric package.

The present note proves that it does. Relative to the already fixed observer-relevant pre-core exact-shell transport core, the only independent sub-pre-core datum still carried by the current image core is the canonical transport section from the recorded pre-core transport body into its sub-pre-core minimizer image. The transport-image body is exactly the image of that section, the restricted projection is its inverse, and the exact-shell transport image is exactly the shell restriction of the same section. Separate listing of those image-core constituents therefore does not add independent benchmark-facing freedom.

The benchmark boundary remains fixed throughout. The one operative metric is still exactly \(\CompletedMetric\), the ordinary matter route is unchanged, there is no second operative metric, the low-energy matter law is not enlarged, and no stronger promotion language is licensed. The positive result is therefore a genuine smaller collapse strictly inside the current image core: the live burden now contracts to an observer-relevant sub-pre-core exact-shell transport-section core. What remains open is to derive or justify that section core from still more primitive substrate structure, or else prove a still sharper collapse strictly inside it.
\end{abstract}

\section{Introduction}

The active derivational program of \TheoryName{} is again constrained by a routing safeguard. Once a current benchmark-facing datum has been compressed to a smaller observer-relevant core, another theorem that merely replays the same datum through a deeper relay does not count as new front-line progress unless it changes benchmark-facing status by more than depth alone. The sub-pre-core image-core theorem made exactly such a compression: it removed the full exact-shell-transport-core-generating sub-pre-core package as the live burden and replaced it by a smaller observer-relevant sub-pre-core exact-shell transport-image core.

That theorem also exposes the next honest question. Inside the image core itself, what is genuinely independent observer-relevant data, and what is merely presentational bookkeeping once the recorded pre-core exact-shell transport core is already fixed? The present note takes the sharper collapse route rather than a same-output deeper-origin relay. It shows that the image-core body, shell image, and restricted projection are not independent pieces at current scope. They are all recovered from a single canonical transport section over the already fixed pre-core transport body.

\paragraph{Framework and claim boundary.}
\TheoryProgram{} denotes the broader \Aether{} / \Aflow{} framework built on the claims that the \Aether{} is the underlying four-dimensional substrate of reality, the \Aflow{} is its intrinsic ordered motion, observed three-dimensional space is the local experiential slice of that deeper substrate, S-time is the experienced order of change arising from matter, light, and the \Aflow{}, and observed expansion is the three-dimensional appearance of deeper four-dimensional ordered motion. Within that framework, \TheoryName{} denotes the exact-closure theory in which the effective gravitational dynamics are adopted to be exactly Einsteinian with universal matter coupling. In the active sequence, \emph{adoption} means use of that established relativistic dynamics without claiming substrate derivation, while \emph{derivation} is reserved for a first-principles recovery from explicit substrate variables.

The present note stays strictly inside that boundary. It does not introduce a second operative metric, it does not enlarge the low-energy matter law, and it does not reopen the already-screened larger sub-pre-core or pre-core packages as if they were again the live burden. Its narrower benchmark-facing role is to isolate the only remaining observer-relevant datum still carried by the current sub-pre-core image core.

\section{Fixed Inputs}

\begin{definition}[Recorded primitive continuation data]
\label{def:fixed}
Fix the following continuation data from the active primitive-reservoir observer-reduction line.
\begin{enumerate}
    \item For each sector label \(\sigma\in\{\Car,\Cur,\Hom\}\), a primitive forcing shell
    \[
    \ForSpace{\sigma},
    \]
    a visible reduction map
    \[
    \Reduction{\sigma}:\ForSpace{\sigma}\to X_{\sigma},
    \]
    and shell-level current and equilibrium carriers
    \[
    \CurrentCarrier{\sigma}:\ForSpace{\sigma}\to Z_{\sigma}^{\mathrm{cur}},
    \qquad
    \EqCarrier{\sigma}:\ForSpace{\sigma}\to Z_{\sigma}^{\mathrm{eq}}.
    \]
    \item The product shell
    \[
    \PrimShell
    \equiv
    \ForSpace{\Car}\times \ForSpace{\Cur}\times \ForSpace{\Hom}.
    \]
    \item The recorded metric and ordinary-matter outputs
    \[
    \MetricOut:\PrimShell\to\{\text{observer metrics}\},
    \qquad
    \MatterOut:\PrimShell\times\MatterSpace\to\{\text{observer stress tensors}\},
    \]
    satisfying
    \begin{equation}
    \MetricOut(\mathbf w)=\CompletedMetric,
    \qquad
    \MatterOut(\mathbf w,\psi)
    =
    T_{\mu\nu}^{\mathrm{matter}}[\CompletedMetric,\psi]
    \label{eq:fixedoutputs}
    \end{equation}
    for every \(\mathbf w\in\PrimShell\) and every matter configuration \(\psi\in\MatterSpace\).
\end{enumerate}
\end{definition}

\begin{definition}[Recorded proto-germ exact-shell target]
\label{def:proto}
Fix \(\sigma\in\{\Car,\Cur,\Hom\}\). The active line already fixes:
\begin{enumerate}
    \item the same substrate target and pairing
    \[
    \SubTarget{\sigma},
    \qquad
    \SubPair{\sigma}:
    \SubTarget{\sigma}\times\SubTarget{\sigma}\to\mathbb R;
    \]
    \item a hidden fiber space \(\HiddenSpace{\sigma}\) with compact hidden body \(\HiddenBody{\sigma}\subset\HiddenSpace{\sigma}\);
    \item a finite-dimensional proto-germ state space \(\ProtoSpace{\sigma}\), a compact convex proto-germ body \(\ProtoBody{\sigma}\subset\ProtoSpace{\sigma}\), a smooth proto-germ restriction section
    \[
    \ProtoSection{\sigma}:\ProtoSpace{\sigma}\to\SubTarget{\sigma},
    \]
    and the exact proto-germ shell
    \[
    \ProtoShellSet{\sigma}
    \equiv
    \bigl(\ProtoSection{\sigma}\bigr)^{-1}(0)\cap\ProtoBody{\sigma};
    \]
    \item a closed embedded proto zero core \(\ProtoZeroCore{\sigma}\subset\ProtoShellSet{\sigma}\) and a split diffeomorphism
    \[
    \ProtoSplit{\sigma}:
    \ProtoZeroCore{\sigma}\times\HiddenBody{\sigma}
    \xrightarrow{\cong}
    \ProtoShellSet{\sigma};
    \]
    \item a hidden charge
    \[
    \ProtoCharge{\sigma}:\HiddenSpace{\sigma}\to\EqnTarget{\sigma},
    \]
    a smooth observer-compatible exact proto-germ zero-response realization
    \[
    \ProtoEmbed{\sigma}:\ForSpace{\sigma}\hookrightarrow\ProtoZeroCore{\sigma},
    \]
    and a fixed-data solvability / coercive package on \(\ProtoZeroCore{\sigma}\), recorded through \(\ProtoKernelCh{\sigma}\) and \(\ProtoStateComp{\sigma}\).
\end{enumerate}
\end{definition}

\begin{definition}[Recorded current transport nucleus]
\label{def:nucleus}
Fix \(\sigma\in\{\Car,\Cur,\Hom\}\). The recorded current transport nucleus \(\PreProtoNucleus\) for sector \(\sigma\) consists of:
\begin{enumerate}
    \item the pre-proto-germ state space \(\PreProtoSpace{\sigma}\), compact convex body \(\PreProtoBody{\sigma}\subset\PreProtoSpace{\sigma}\), projection
    \[
    \PreProtoProj{\sigma}:\PreProtoSpace{\sigma}\to\ProtoSpace{\sigma},
    \]
    and proto section
    \[
    \PreProtoProtoMin{\sigma}:\ProtoBody{\sigma}\to\PreProtoBody{\sigma}
    \]
    satisfying
    \[
    \PreProtoProj{\sigma}\bigl(\PreProtoBody{\sigma}\bigr)=\ProtoBody{\sigma},
    \qquad
    \PreProtoProj{\sigma}\circ\PreProtoProtoMin{\sigma}
    =
    \mathrm{id}_{\ProtoBody{\sigma}};
    \]
    \item the restriction section
    \[
    \PreProtoSection{\sigma}:\PreProtoSpace{\sigma}\to\SubTarget{\sigma},
    \]
    the induced functional
    \[
    \PreProtoFunctional{\sigma}(q)
    \equiv
    \SubPair{\sigma}\bigl(\PreProtoSection{\sigma}(q),\PreProtoSection{\sigma}(q)\bigr),
    \]
    and the exact shell
    \[
    \PreProtoShellSet{\sigma}
    \equiv
    \bigl(\PreProtoSection{\sigma}\bigr)^{-1}(0)\cap\PreProtoBody{\sigma},
    \]
    with
    \[
    \PreProtoSection{\sigma}\circ\PreProtoProtoMin{\sigma}
    =
    \ProtoSection{\sigma},
    \qquad
    \PreProtoProj{\sigma}\big|_{\PreProtoShellSet{\sigma}}
    :
    \PreProtoShellSet{\sigma}
    \xrightarrow{\cong}
    \ProtoShellSet{\sigma};
    \]
    \item benchmark faithfulness, meaning preservation of the benchmark outputs \eqref{eq:fixedoutputs}, no second operative metric, no enlarged low-energy matter law, and no premature promotion from adoption to completed derivation.
\end{enumerate}
\end{definition}

\begin{definition}[Recorded observer-relevant pre-nucleus exact-shell core]
\label{def:core}
Fix \(\sigma\in\{\Car,\Cur,\Hom\}\). The observer-relevant pre-nucleus exact-shell core \(\PreNucCore\) for sector \(\sigma\) consists of:
\begin{enumerate}
    \item the pre-nucleus state space \(\PreNucSpace{\sigma}\), compact convex body \(\PreNucBody{\sigma}\subset\PreNucSpace{\sigma}\), affine surjection
    \[
    \PreNucProj{\sigma}:\PreNucSpace{\sigma}\to\PreProtoSpace{\sigma},
    \]
    and fiber minimizer
    \[
    \PreNucNucleusMin{\sigma}:\PreProtoBody{\sigma}\to\PreNucBody{\sigma}
    \]
    satisfying
    \[
    \PreNucProj{\sigma}\bigl(\PreNucBody{\sigma}\bigr)=\PreProtoBody{\sigma},
    \qquad
    \PreNucProj{\sigma}\circ\PreNucNucleusMin{\sigma}
    =
    \mathrm{id}_{\PreProtoBody{\sigma}};
    \]
    \item a closed embedded exact pre-nucleus shell
    \[
    \PreNucShellSet{\sigma}\subset\PreNucBody{\sigma}
    \]
    satisfying
    \[
    \PreNucNucleusMin{\sigma}\bigl(\PreProtoShellSet{\sigma}\bigr)
    \subset
    \PreNucShellSet{\sigma},
    \qquad
    \PreNucProj{\sigma}\big|_{\PreNucShellSet{\sigma}}
    :
    \PreNucShellSet{\sigma}
    \xrightarrow{\cong}
    \PreProtoShellSet{\sigma};
    \]
    \item benchmark faithfulness in the sense of Definition \ref{def:nucleus}.
\end{enumerate}
\end{definition}

\begin{definition}[Recorded pre-core exact-shell transport core]
\label{def:transportcore}
Fix \(\sigma\in\{\Car,\Cur,\Hom\}\). The recorded observer-relevant pre-core exact-shell transport core \(\PreCoreTransportCore\) for sector \(\sigma\) consists of:
\begin{enumerate}
    \item a compact transport-core body
    \[
    \PreCoreTrBody{\sigma},
    \]
    an affine projection
    \[
    \PreCoreTrProj{\sigma}:\PreCoreTrBody{\sigma}\to\PreNucBody{\sigma},
    \]
    and a smooth transport-core lift
    \[
    \PreCoreTrLift{\sigma}:\PreNucBody{\sigma}\to\PreCoreTrBody{\sigma}
    \]
    satisfying
    \[
    \PreCoreTrProj{\sigma}\bigl(\PreCoreTrBody{\sigma}\bigr)=\PreNucBody{\sigma},
    \qquad
    \PreCoreTrProj{\sigma}\circ\PreCoreTrLift{\sigma}
    =
    \mathrm{id}_{\PreNucBody{\sigma}};
    \]
    \item a closed embedded exact transport shell
    \[
    \PreCoreTrShell{\sigma}\subset\PreCoreTrBody{\sigma}
    \]
    satisfying
    \[
    \PreCoreTrLift{\sigma}\bigl(\PreNucShellSet{\sigma}\bigr)
    \subset
    \PreCoreTrShell{\sigma},
    \qquad
    \PreCoreTrProj{\sigma}\big|_{\PreCoreTrShell{\sigma}}
    :
    \PreCoreTrShell{\sigma}
    \xrightarrow{\cong}
    \PreNucShellSet{\sigma};
    \]
    \item benchmark faithfulness in the sense of Definition \ref{def:core}.
\end{enumerate}
\end{definition}

\begin{remark}
Definitions \ref{def:fixed}--\ref{def:transportcore} are fixed inputs. The one operative metric, the ordinary matter route, the fixed proto-germ exact-shell target, the current transport nucleus, the recorded pre-nucleus exact-shell core, and the recorded pre-core exact-shell transport core are not reopened here. The present note asks only which part of the current sub-pre-core image-core presentation still carries independent benchmark-facing burden once those downstream data are already fixed.
\end{remark}

\section{Full Sub-Pre-Core Package}

\begin{definition}[Exact-shell-transport-core-generating substrate sub-pre-core dynamics]
\label{def:subpre}
Fix \(\sigma\in\{\Car,\Cur,\Hom\}\). An exact-shell-transport-core-generating substrate sub-pre-core dynamics package \(\SubPreCorePackage\) for sector \(\sigma\) consists of the following data.
\begin{enumerate}
    \item A finite-dimensional Euclidean sub-pre-core state space \(\SubPreCoreSpace{\sigma}\), a compact convex sub-pre-core body
    \[
    \SubPreCoreBody{\sigma}\subset\SubPreCoreSpace{\sigma}
    \]
    with nonempty interior, an affine surjection
    \[
    \SubPreCoreProj{\sigma}:\SubPreCoreSpace{\sigma}\to\PreCoreTrBody{\sigma},
    \]
    and a smooth sub-pre-core minimizer
    \[
    \SubPreCoreTransportMin{\sigma}:\PreCoreTrBody{\sigma}\to\SubPreCoreBody{\sigma}
    \]
    satisfying
    \[
    \SubPreCoreProj{\sigma}\bigl(\SubPreCoreBody{\sigma}\bigr)=\PreCoreTrBody{\sigma},
    \qquad
    \SubPreCoreProj{\sigma}\circ\SubPreCoreTransportMin{\sigma}
    =
    \mathrm{id}_{\PreCoreTrBody{\sigma}}.
    \]
    \item A closed embedded exact sub-pre-core shell
    \[
    \SubPreCoreShellSet{\sigma}\subset\SubPreCoreBody{\sigma}
    \]
    satisfying
    \[
    \SubPreCoreTransportMin{\sigma}\bigl(\PreCoreTrShell{\sigma}\bigr)
    \subset
    \SubPreCoreShellSet{\sigma},
    \qquad
    \SubPreCoreProj{\sigma}\big|_{\SubPreCoreShellSet{\sigma}}
    :
    \SubPreCoreShellSet{\sigma}
    \xrightarrow{\cong}
    \PreCoreTrShell{\sigma}.
    \]
    \item Benchmark faithfulness, meaning preservation of the benchmark outputs \eqref{eq:fixedoutputs}, no second operative metric, no enlarged low-energy matter law, and no premature promotion from adoption to completed derivation.
\end{enumerate}
\end{definition}

\begin{remark}
Definition \ref{def:subpre} is the full current sub-pre-core package. The previous active-primary theorem already showed that the full package is not the minimal benchmark-facing datum once the recorded pre-core transport core is fixed.
\end{remark}

\section{Observer-Relevant Sub-Pre-Core Exact-Shell Transport-Image Core}

\begin{definition}[Observer-relevant sub-pre-core exact-shell transport-image core]
\label{def:imagecore}
Fix \(\sigma\in\{\Car,\Cur,\Hom\}\) and a benchmark-faithful sub-pre-core package in the sense of Definition \ref{def:subpre}. The \emph{observer-relevant sub-pre-core exact-shell transport-image core} \(\SubPreImageCore\) for sector \(\sigma\) consists of:
\begin{enumerate}
    \item the compact transport-image body
    \[
    \SubPreImgBody{\sigma}
    \equiv
    \SubPreCoreTransportMin{\sigma}\bigl(\PreCoreTrBody{\sigma}\bigr)
    \subset
    \SubPreCoreBody{\sigma},
    \]
    together with the restricted projection
    \[
    \SubPreImgProj{\sigma}
    \equiv
    \SubPreCoreProj{\sigma}\big|_{\SubPreImgBody{\sigma}}
    :
    \SubPreImgBody{\sigma}\to\PreCoreTrBody{\sigma};
    \]
    \item the exact-shell transport image
    \[
    \SubPreImgShell{\sigma}
    \equiv
    \SubPreCoreTransportMin{\sigma}\bigl(\PreCoreTrShell{\sigma}\bigr)
    \subset
    \SubPreCoreShellSet{\sigma},
    \]
    with the restricted shell projection
    \[
    \SubPreImgProj{\sigma}\big|_{\SubPreImgShell{\sigma}}
    :
    \SubPreImgShell{\sigma}\to\PreCoreTrShell{\sigma};
    \]
    \item benchmark faithfulness in the sense of Definition \ref{def:subpre}.
\end{enumerate}
\end{definition}

\begin{remark}
Definition \ref{def:imagecore} is already smaller than the full sub-pre-core package. The present note asks whether it is itself still larger than necessary once the recorded pre-core exact-shell transport body and shell are treated as fixed upstream data.
\end{remark}

\section{Observer-Relevant Sub-Pre-Core Exact-Shell Transport-Section Core}

\begin{definition}[Observer-relevant sub-pre-core exact-shell transport-section core]
\label{def:sectioncore}
Fix \(\sigma\in\{\Car,\Cur,\Hom\}\) and a benchmark-faithful sub-pre-core package in the sense of Definition \ref{def:subpre}. The \emph{observer-relevant sub-pre-core exact-shell transport-section core} \(\SubPreSectionCore\) for sector \(\sigma\) consists of:
\begin{enumerate}
    \item the canonical transport section
    \[
    \SubPreSection{\sigma}
    \equiv
    \SubPreCoreTransportMin{\sigma}
    :
    \PreCoreTrBody{\sigma}\to\SubPreCoreBody{\sigma};
    \]
    \item benchmark faithfulness in the sense of Definition \ref{def:subpre}.
\end{enumerate}
The shell restriction
\[
\SubPreSection{\sigma}\big|_{\PreCoreTrShell{\sigma}}
:
\PreCoreTrShell{\sigma}\to\SubPreCoreShellSet{\sigma}
\]
is understood as derived data rather than separately listed core data.
\end{definition}

\begin{remark}
Definition \ref{def:sectioncore} is strictly smaller than Definition \ref{def:imagecore} at current scope. Once the recorded transport-core body \(\PreCoreTrBody{\sigma}\) and shell \(\PreCoreTrShell{\sigma}\) are fixed, the image body, image shell, and restricted projection are recovered from the single canonical section \(\SubPreSection{\sigma}\); they need not be listed independently.
\end{remark}

\section{Collapse of the Image Core}

\begin{proposition}[The canonical transport section reconstructs the image body and projection]
\label{prop:reconstructbody}
For every benchmark-faithful sub-pre-core package,
\[
\SubPreImgBody{\sigma}
=
\SubPreSection{\sigma}\bigl(\PreCoreTrBody{\sigma}\bigr),
\]
and the restricted projection
\[
\SubPreImgProj{\sigma}:\SubPreImgBody{\sigma}\to\PreCoreTrBody{\sigma}
\]
is a homeomorphism with inverse \(\SubPreSection{\sigma}\).
\end{proposition}

\begin{proof}
The displayed identity is Definition \ref{def:imagecore}(1) together with Definition \ref{def:sectioncore}(1). By Definition \ref{def:subpre}(1),
\[
\SubPreCoreProj{\sigma}\circ\SubPreSection{\sigma}
=
\mathrm{id}_{\PreCoreTrBody{\sigma}}.
\]
Restricting \(\SubPreCoreProj{\sigma}\) to \(\SubPreImgBody{\sigma}\) yields
\[
\SubPreImgProj{\sigma}\circ\SubPreSection{\sigma}
=
\mathrm{id}_{\PreCoreTrBody{\sigma}}.
\]
Conversely, every \(y\in\SubPreImgBody{\sigma}\) has the form
\[
y=\SubPreSection{\sigma}(w)
\qquad
\text{for some }
w\in\PreCoreTrBody{\sigma}.
\]
Hence
\begin{align*}
\SubPreSection{\sigma}\bigl(\SubPreImgProj{\sigma}(y)\bigr)
&=
\SubPreSection{\sigma}\bigl(\SubPreCoreProj{\sigma}(y)\bigr)\\
&=
\SubPreSection{\sigma}\bigl(\SubPreCoreProj{\sigma}(\SubPreSection{\sigma}(w))\bigr)\\
&=
\SubPreSection{\sigma}(w)\\
&=
y.
\end{align*}
Therefore \(\SubPreImgProj{\sigma}\) and \(\SubPreSection{\sigma}\) are mutually inverse continuous maps, so the restricted projection is a homeomorphism with inverse \(\SubPreSection{\sigma}\).
\end{proof}

\begin{proposition}[The shell image is exactly the shell restriction of the canonical section]
\label{prop:reconstructshell}
For every benchmark-faithful sub-pre-core package,
\[
\SubPreImgShell{\sigma}
=
\SubPreSection{\sigma}\bigl(\PreCoreTrShell{\sigma}\bigr),
\]
and the shell restriction
\[
\SubPreSection{\sigma}\big|_{\PreCoreTrShell{\sigma}}
:
\PreCoreTrShell{\sigma}\xrightarrow{\cong}\SubPreImgShell{\sigma}
\]
is a diffeomorphism with inverse
\[
\SubPreImgProj{\sigma}\big|_{\SubPreImgShell{\sigma}}
:
\SubPreImgShell{\sigma}\xrightarrow{\cong}\PreCoreTrShell{\sigma}.
\]
\end{proposition}

\begin{proof}
The displayed equality is Definition \ref{def:imagecore}(2) together with Definition \ref{def:sectioncore}(1). By Definition \ref{def:subpre}(2),
\[
\SubPreCoreProj{\sigma}\big|_{\SubPreCoreShellSet{\sigma}}
:
\SubPreCoreShellSet{\sigma}\xrightarrow{\cong}\PreCoreTrShell{\sigma}
\]
is a diffeomorphism, and
\[
\SubPreSection{\sigma}\bigl(\PreCoreTrShell{\sigma}\bigr)
\subset
\SubPreCoreShellSet{\sigma}.
\]
Restricting the shell diffeomorphism to \(\SubPreImgShell{\sigma}\) gives
\[
\SubPreImgProj{\sigma}\big|_{\SubPreImgShell{\sigma}}
\circ
\SubPreSection{\sigma}\big|_{\PreCoreTrShell{\sigma}}
=
\mathrm{id}_{\PreCoreTrShell{\sigma}}.
\]
Conversely, if \(z\in\SubPreImgShell{\sigma}\), then \(z=\SubPreSection{\sigma}(w)\) for some \(w\in\PreCoreTrShell{\sigma}\), so
\[
\SubPreSection{\sigma}\bigl(\SubPreImgProj{\sigma}(z)\bigr)
=
\SubPreSection{\sigma}(w)
=
z.
\]
Thus the shell restriction of the canonical section and the restricted shell projection are inverse diffeomorphisms.
\end{proof}

\begin{proposition}[The full image core is reconstructed from the section core]
\label{prop:reconstructimage}
For every benchmark-faithful sub-pre-core package, the full observer-relevant sub-pre-core exact-shell transport-image core \(\SubPreImageCore\) is reconstructed from the smaller section core \(\SubPreSectionCore\) together with the fixed transport-core data \(\PreCoreTrBody{\sigma}\) and \(\PreCoreTrShell{\sigma}\).
\end{proposition}

\begin{proof}
By Proposition \ref{prop:reconstructbody}, the image body is the image of the canonical section and the restricted projection is its inverse. By Proposition \ref{prop:reconstructshell}, the shell image is the shell restriction of that same section. Benchmark faithfulness is already part of Definition \ref{def:sectioncore}(2). Therefore each constituent of Definition \ref{def:imagecore} is recovered from Definition \ref{def:sectioncore} plus the fixed upstream transport-core body and shell.
\end{proof}

\begin{proposition}[Downstream exact-shell transport factors through the section core]
\label{prop:downstream}
For every benchmark-faithful sub-pre-core package, the recovery of the recorded pre-core exact-shell transport core and all downstream shell transport to the fixed pre-nucleus exact-shell core, transport nucleus, and proto-germ exact-shell target are already carried by the section core \(\SubPreSectionCore\) together with the fixed transport-core data.
\end{proposition}

\begin{proof}
Proposition \ref{prop:reconstructimage} reconstructs the full image core from the section core. The recorded transport-core projection
\[
\PreCoreTrProj{\sigma}:\PreCoreTrBody{\sigma}\to\PreNucBody{\sigma}
\]
and lift
\[
\PreCoreTrLift{\sigma}:\PreNucBody{\sigma}\to\PreCoreTrBody{\sigma}
\]
are already fixed by Definition \ref{def:transportcore}. Hence the composite
\[
\PreCoreTrProj{\sigma}\circ\SubPreImgProj{\sigma}
:
\SubPreImgBody{\sigma}\to\PreNucBody{\sigma}
\]
and its shell restriction are recovered from the section core because \(\SubPreImgProj{\sigma}=(\SubPreSection{\sigma})^{-1}\) on \(\SubPreImgBody{\sigma}\). The same argument then composes with the fixed shell diffeomorphisms of Definitions \ref{def:core} and \ref{def:nucleus} to recover downstream transport to \(\PreProtoShellSet{\sigma}\) and \(\ProtoShellSet{\sigma}\). Therefore nothing benchmark-facing downstream of the image core requires more than the section core together with the already recorded transport-core package.
\end{proof}

\begin{proposition}[Separate image-core presentation is benchmark neutral at current scope]
\label{prop:neutral}
For every benchmark-faithful sub-pre-core package, the separate presentation of
\[
\SubPreImgBody{\sigma},
\qquad
\SubPreImgProj{\sigma},
\qquad
\SubPreImgShell{\sigma}
\]
does not add independent benchmark-facing freedom beyond the section core \(\SubPreSectionCore\).
\end{proposition}

\begin{proof}
By Proposition \ref{prop:reconstructimage}, the listed constituents are reconstructed from the canonical transport section together with the already fixed transport-core body and shell. By Proposition \ref{prop:downstream}, all downstream transport relevant to the recorded pre-nucleus exact-shell core, current transport nucleus, and fixed proto-germ exact-shell target is already carried by that reconstructed data. Therefore the image-core presentation as a separate body, shell, and projection is not additional benchmark-facing information beyond the smaller section core.
\end{proof}

\begin{proposition}[The benchmark package remains fixed]
\label{prop:benchmark}
Every observer-relevant sub-pre-core exact-shell transport-section core preserves the same benchmark one-metric package \eqref{eq:fixedoutputs}. In particular, the operative metric remains exactly \(\CompletedMetric\), the ordinary matter route is unchanged, there is no second operative metric, and the low-energy matter law is not enlarged.
\end{proposition}

\begin{proof}
This is part of benchmark faithfulness in Definition \ref{def:sectioncore}(2). The present note only removes presentational redundancy inside the current image core. It does not alter the benchmark exact-closure package.
\end{proof}

\section{Main Theorem}

\begin{theorem}[Observer-relevant sub-pre-core exact-shell transport-image core collapse to an observer-relevant sub-pre-core exact-shell transport-section core]
\label{thm:main}
Fix the primitive continuation data of Definition \ref{def:fixed}, the recorded proto-germ exact-shell target of Definition \ref{def:proto}, the recorded current transport nucleus of Definition \ref{def:nucleus}, the recorded observer-relevant pre-nucleus exact-shell core of Definition \ref{def:core}, the recorded observer-relevant pre-core exact-shell transport core of Definition \ref{def:transportcore}, and a benchmark-faithful sub-pre-core package in the sense of Definition \ref{def:subpre}. Then, for each sector \(\sigma\in\{\Car,\Cur,\Hom\}\), the live benchmark-facing burden inside the current sub-pre-core image-core frontier collapses from the observer-relevant sub-pre-core exact-shell transport-image core \(\SubPreImageCore\) to the smaller observer-relevant sub-pre-core exact-shell transport-section core \(\SubPreSectionCore\). More precisely:
\begin{enumerate}
    \item the image body and restricted projection are reconstructed from the canonical transport section by Proposition \ref{prop:reconstructbody};
    \item the exact-shell transport image is exactly the shell restriction of that section by Proposition \ref{prop:reconstructshell};
    \item the full image core is therefore recovered from the section core and the fixed transport-core body and shell by Proposition \ref{prop:reconstructimage};
    \item all downstream recovery of the recorded pre-core exact-shell transport core, the pre-nucleus exact-shell core, the current transport nucleus, and the fixed proto-germ exact-shell target factors through the section core by Proposition \ref{prop:downstream};
    \item the separate image-core presentation adds no independent benchmark-facing freedom beyond the section core by Proposition \ref{prop:neutral};
    \item the benchmark boundary remains fixed by Proposition \ref{prop:benchmark};
    \item consequently the remaining honest burden below the previous image-core frontier is no longer the full observer-relevant sub-pre-core exact-shell transport-image core \(\SubPreImageCore\) as such. What remains is to derive or justify the sharper section core \(\SubPreSectionCore\) from still more primitive substrate structure, or else prove a still smaller theorem-level collapse strictly inside that section core.
\end{enumerate}
\end{theorem}

\begin{proof}
Item (1) is Proposition \ref{prop:reconstructbody}. Item (2) is Proposition \ref{prop:reconstructshell}. Item (3) is Proposition \ref{prop:reconstructimage}. Item (4) is Proposition \ref{prop:downstream}. Item (5) is Proposition \ref{prop:neutral}. Item (6) is Proposition \ref{prop:benchmark}. Item (7) is the routing consequence of the previous items together with the active safeguard against counting another same-output deeper-origin relay as new front-line progress when the live benchmark-relevant datum is unchanged.
\end{proof}

\begin{corollary}[What must be done next]
\label{cor:next}
After Theorem \ref{thm:main}, the next honest benchmark-facing manuscript must either:
\begin{enumerate}
    \item derive or justify the observer-relevant sub-pre-core exact-shell transport-section core \(\SubPreSectionCore\) from still more primitive substrate structure in a way that changes benchmark-facing status;
    \item identify a genuinely smaller theorem-level observer-relevant datum strictly inside \(\SubPreSectionCore\) and prove a sharper collapse to it;
    \item do either of the previous two while preserving the same benchmark one-metric package and the same claim boundary.
\end{enumerate}
Another same-output deeper-origin relay that preserves this section core, another re-expansion back to the larger image-core presentation as if its recovered constituents were still independently live, another reopening of the already-screened full sub-pre-core or pre-core packages as if they were again the active burden, or any move that introduces a second operative metric, enlarges the ordinary matter law, or uses stronger promotion language does not satisfy the present burden. If a quantum continuation is pursued, it matters benchmark-facingly only insofar as it derives this same section core, collapses it further, or closes a benchmark derivation gate in a genuinely new way.
\end{corollary}

\begin{proof}
Theorem \ref{thm:main} shows that the image-core presentation is no longer the minimal benchmark-facing datum at this stage. Therefore a subsequent manuscript must improve on the smaller section core rather than merely restating the larger image core.
\end{proof}

\section{Conclusion}

The positive result of this note is a disciplined second compression step inside the current sub-pre-core frontier. The previous active-primary theorem already showed that the full sub-pre-core package is not the minimal benchmark-facing datum, because its ambient presentation collapses to an observer-relevant exact-shell transport-image core. The present theorem then takes the honest next internal route available there: it shows that even that image-core presentation is still larger than necessary.

What survives as the live burden is the observer-relevant sub-pre-core exact-shell transport-section core. Relative to the already fixed pre-core exact-shell transport body and shell, that sharper datum is just the canonical transport section that picks out the sub-pre-core minimizer image. The image body is its image, the exact-shell transport image is its shell restriction, and the restricted projection is simply its inverse. Those separately listed image-core constituents therefore do not add independent benchmark-facing freedom.

This is the honest next burden after the present theorem. A new primary manuscript must derive or justify that sharper section core from still more primitive substrate structure, or else collapse it further, while preserving the same benchmark one-metric package and the same claim boundary. The current result does not yet complete a first-principles derivation of GR from substrate microphysics, but it does narrow the live derivational burden one level further on the active line.

\end{document}
